% Generated draft: requires mathematical and source review.
\documentclass[11pt]{article}
\usepackage[margin=25mm]{geometry}
\usepackage{fontspec}
\setmainfont{Latin Modern Roman}
\usepackage{amsmath,amssymb,mathtools,unicode-math}
\setmathfont{Latin Modern Math}
\usepackage{enumitem,needspace,xurl,hyperref}
\setlength{\parindent}{0pt}
\setlength{\parskip}{5pt}
\setlength{\emergencystretch}{3em}
\newcommand{\N}{\mathbb{N}}
\newcommand{\Z}{\mathbb{Z}}
\newcommand{\Q}{\mathbb{Q}}
\newcommand{\R}{\mathbb{R}}
\newcommand{\C}{\mathbb{C}}
\newcommand{\F}{\mathbb{F}}
\newcommand{\A}{\mathbb{A}}
\newcommand{\PP}{\mathbb{P}}
\newcommand{\E}{\mathbb{E}}
\newcommand{\eps}{\varepsilon}
\newcommand{\dd}{\,\mathrm{d}}
\begin{document}
\section*{UnsolvedMath ID: 30006344}
\textbf{Source problem code:} catalogue entry
\subsection{Definitions and mathematical statement}
Let \(G\) be a finitely presented group and let \(R=\Z[G]\) be its integral group ring. If \(X\) is a finite connected two-dimensional CW-complex with a specified isomorphism \(\pi_1(X)\cong G\), its universal cover \(\widetilde X\) has an augmented cellular chain complex
\[
C_2(\widetilde X;\Z)\xrightarrow{d_2}C_1(\widetilde X;\Z)\xrightarrow{d_1}C_0(\widetilde X;\Z)\xrightarrow{\varepsilon}\Z\longrightarrow0.
\]
Here the deck transformations give a left \(R\)-module convention, each \(C_i(\widetilde X;\Z)\) is finitely generated free over \(R\), and \(\Z\) is the trivial \(R\)-module. The complex is exact at \(C_1\), \(C_0\), and \(\Z\); injectivity of \(d_2\) is not required. The cellular complex relevant to the group ring is that of \(\widetilde X\).

For an augmented complex \(A_2\to A_1\to A_0\to\Z\to0\), literal cellular realization means that there exist such an \(X\), a marking \(\pi_1(X)\cong G\), and an isomorphism of augmented \(R\)-chain complexes from this complex to \(C_*(\widetilde X;\Z)\to\Z\), inducing the identity on the augmentation target.

\textbf{Problem, with unresolved source qualifications.} For every finitely presented group \(G\), characterize which algebraic two-complexes over \(\Z[G]\) admit geometric realization by finite connected two-dimensional CW-complexes with fundamental group identified with \(G\).

The supplied record does not define the admissible algebraic two-complexes: in particular, it does not specify their finiteness, freeness, stable freeness or projectivity requirements, or their augmentation and exactness conditions. It also does not specify whether realization requires an actual chain isomorphism, as in the literal interpretation defined above, or only a chain homotopy equivalence, possibly after stabilization. These distinctions are part of the unresolved formulation; the literal interpretation is not asserted to be the intended version. The original excerpt refers only to \emph{certain} algebraic two-complexes, whereas the curated statement requests a characterization. The omitted restriction must be recovered before this can be treated as a fully specified realization problem.
\subsection{Short English statement}
For a finitely presented group \(G\), determine which algebraic two-complexes over \(\Z[G]\) come from universal covers of finite two-dimensional CW-complexes with fundamental group \(G\). The supplied record leaves both the admissible algebraic complexes and the required notion of equivalence unresolved. Its open-status assessment is supplied dated curation, requiring independent verification.
\subsection{Sources}
\begin{itemize}
\item Local UnsolvedMath registry, supplied complete source record, permanent ID \(30006344\), source code \texttt{catalogue entry}; collection: Oberwolfach Reports Open Problems. Title: \emph{Geometric Realizability of Algebraic Two-Complexes over Group Rings}.
\item Supplied report citation: \emph{Computational Group Theory} (2025), Oberwolfach Reports, DOI \texttt{10.4171/owr/2025/27}, pp.\ 1389--1450. \url{https://doi.org/10.4171/owr/2025/27}. No problem-specific page is supplied.
\item Supplied literature source: \emph{Oberwolfach Report 27/2025}. \url{https://ems.press/content/serial-article-files/51860}.
\item Supplied literature source: \emph{Stable modules and Wall's D(2)-problem}. \url{https://ems.press/journals/cmh/articles/380}.
\item The original excerpt appends the presentation title \emph{Finding exotic presentations of quaternion groups}, followed by Tommy Hofmann, joint work with John Nicholson. The record does not explicitly identify a proposer; this presentation attribution is preserved without assigning authorship of the general realization problem.
\item Literature assessment dated exactly \texttt{2026-08-22}: the supplied curation labels the problem open, identifies it with Wall's algebraic 2-complex realization problem, and reports reductions and special finite-group cases rather than a universal realization theorem. This is supplied dated curation derived from the OWR record, not an independently verified literature assessment; its status, identification, and cited claims require independent checking.
\end{itemize}
\end{document}
